\origpage{272}{ROBERTSON AND SEYMOUR}
\hspace*{12pt}(1) \textit{If $r, r' \in R(K)$ with $r, r' \subseteq \Sigma - Z$, then there is a sequence $r = r_1, r_2, \dots, r_t = r'$ of regions of $K$, all included in $\Sigma - Z$, such that for $1 \le i < t$ there is an edge of $K$ incident with $r_i$ and $r_{i+1}$.}

Now for each $e \in E(\Gamma)$ there is a region $r$ of $K$ with $e \subseteq r$; and $e \in E(\Gamma')$ if and only if $r \not\subseteq Z$, that is, $r \subseteq \Sigma - Z$. Conversely, every region of $K$ in $\Sigma - Z$ includes an edge of $\Gamma'$. From (1), since $\Sigma - Z$ is open and non-null, we deduce that $\Gamma'$ is non-null and connected. To show that $\Gamma'$ is 2-cell, it therefore suffices to show that if $\phi \colon S^1 \to \Sigma$ is continuous and $|\phi| \cap U(\Gamma') = \emptyset$ then $\phi$ is null-homotopic in $\Sigma$. But it is easily seen that $\phi$ is homotopic in $\Sigma$ to a continuous $\phi' \colon S^1 \to \Sigma$ with $|\phi'| \subseteq Z$; and since $Z$ is simply-connected, $\phi'$ (and hence $\phi$) is null-homotopic in $\Sigma$. The result follows. $\square$

\sectionhead{\textsc{References}}

1. D. B. A. Epstein, Curves on 2-manifolds and isotopies, \textit{Acta Math.} \textbf{115} (1966), 83-107. 2. N. Robertson and P. D. Seymour, Graph minors. VI. Disjoint paths across a disc, \textit{J. Combin. Theory Ser. B} \textbf{41} (1986), 115-138. 3. N. Robertson and P. D. Seymour, Graph minors. VII. Disjoint paths on a surface, \textit{J. Combin. Theory Ser. B} \textbf{45} (1988), 212-254. 4. N. Robertson and P. D. Seymour, Graph minors. X. Obstructions to tree-decomposition, \textit{J. Combin. Theory Ser. B} \textbf{52} (1991), 153-190. 5. N. Robertson and P. D. Seymour, Graph minors. XI. Circuits on a surface, \textit{J. Combin. Theory Ser. B} \textbf{60} (1994), 72-106.
